\documentclass[11pt]{article}
\usepackage[margin=1in]{geometry}
\usepackage{xcolor}
\usepackage{tcolorbox}
\usepackage{hyperref}
\usepackage{enumitem}
\setlist{noitemsep,topsep=0pt}
\usepackage{booktabs}
\usepackage{array}
\usepackage{listings}
\usepackage{amsmath}

% Define OSU orange color
\definecolor{osaorange}{RGB}{220,115,38}

% Configure hyperref colors
\hypersetup{
    colorlinks=true,
    linkcolor=blue,
    urlcolor=blue,
    citecolor=blue
}

% Configure R code listings
\lstset{
    language=R,
    basicstyle=\ttfamily\small,
    breaklines=true,
    breakatwhitespace=true,
    postbreak=\mbox{\textcolor{red}{$\hookrightarrow$}\space},
    frame=single,
    backgroundcolor=\color{gray!10},
    numbers=left,
    numberstyle=\tiny,
    keywordstyle=\color{blue},
    commentstyle=\color{green!60!black},
    stringstyle=\color{red},
    columns=flexible,
    keepspaces=true,
    xleftmargin=0pt,
    xrightmargin=0pt,
    linewidth=\textwidth,
    aboveskip=\baselineskip,
    belowskip=\baselineskip,
    framexleftmargin=0pt,
    framexrightmargin=0pt
}

\title{}
\author{}
\date{}

\begin{document}

% Orange header box with black outline
\begin{tcolorbox}[colback=osaorange,colframe=black,boxrule=2pt,arc=0pt,left=6pt,right=6pt,top=10pt,bottom=10pt]
\centering
\textcolor{white}{\textbf{\Large In-Class Worksheet}}\\
\textcolor{white}{\textbf{\large Introduction to R and Data Presentation}}
\end{tcolorbox}

\vspace{8pt}

\noindent\textbf{Name:} \rule{8cm}{0.4pt} \hfill \textbf{Date:} \rule{4cm}{0.4pt}\\[0.5ex]
\textbf{Course:} H524 Introduction to Biostatistics\\
\textbf{Topic:} Introduction to R, Data Presentation, Numerical Summary Measures

\vspace{8pt}

\section{Part 1: Getting Started with R}

\subsection{Exercise 1.1: Basic R Operations}
Work with a partner to complete these basic R calculations:

\begin{lstlisting}
# 1. Use R as a calculator
3 * 2 + 4

# 2. Create a simple vector of systolic blood pressure readings
blood_pressure <- c(170, 193, 110, 135, 135, 115)
\end{lstlisting}

\textbf{Discussion Questions:}
\begin{itemize}
\item How do you create a vector in R?
\item What does the \texttt{c()} function do?
\end{itemize}

\section{Part 2: Data Types and Classification}

\subsection{Exercise 2.1: Identifying Data Types}
For each variable below, classify it as:
\begin{itemize}
\item \textbf{Nominal} (categorical without order)
\item \textbf{Ordinal} (categorical with order)
\item \textbf{Discrete} (countable numbers)
\item \textbf{Continuous} (measurable numbers)
\end{itemize}

\begin{center}
\begin{tabular}{|>{\raggedright\arraybackslash}p{4.8cm}|>{\centering\arraybackslash}p{2.2cm}|>{\raggedright\arraybackslash}p{5.5cm}|}
\hline
\textbf{Variable} & \textbf{Type} & \textbf{Reasoning} \\
\hline
Blood type (A, B, AB, O) & & \\[2ex]
\hline
Pain level (none, mild, moderate, severe) & & \\[2ex]
\hline
Number of hospitalizations this year & & \\[2ex]
\hline
Body temperature (°F) & & \\[2ex]
\hline
Treatment group (control, low dose, high dose) & & \\[2ex]
\hline
\end{tabular}
\end{center}

\section{Part 3: Descriptive Statistics}

\subsection{Exercise 3.1: Measures of Central Tendency}
Using the blood pressure data: \texttt{blood\_pressure <- c(170, 193, 110, 135, 135, 115)}

Complete the calculations both by hand and in R:

\begin{lstlisting}
# Calculate mean
mean_bp <-

# Calculate median
median_bp <-

# Find the mode (most frequent value)
# Hint: Use table() function

# Order the data
ordered_bp <-
\end{lstlisting}

\textbf{Hand Calculations:}
\begin{itemize}
\item Mean: \rule{3cm}{0.4pt} mmHg
\item Median: \rule{3cm}{0.4pt} mmHg
\item Mode: \rule{3cm}{0.4pt} mmHg
\end{itemize}

\textbf{R Results:}
\begin{itemize}
\item Mean: \rule{3cm}{0.4pt} mmHg
\item Median: \rule{3cm}{0.4pt} mmHg
\item Mode: \rule{3cm}{0.4pt} mmHg
\end{itemize}

\newpage
\section{Part 4: Measures of Variability}
\subsection{Exercise 4.1: Spread of Data}
\begin{lstlisting}
# Calculate range
max_bp <- max(blood_pressure)
min_bp <- min(blood_pressure)
range_bp <- max_bp - min_bp

# Calculate variance
var_bp <-

# Calculate standard deviation
sd_bp <-

# Calculate standard error of the mean
n <- length(blood_pressure)
sem_bp <- sd_bp / sqrt(n)
\end{lstlisting}

\textbf{Fill in your results:}
\begin{itemize}
\item Range: \rule{3cm}{0.4pt} mmHg
\item Variance: \rule{3cm}{0.4pt} mmHg$^2$
\item Standard Deviation: \rule{3cm}{0.4pt} mmHg
\item Standard Error of Mean: \rule{3cm}{0.4pt} mmHg
\end{itemize}

\textbf{Discussion:} What do these measures tell us about the blood pressure data?

\section{Wrap-up Discussion}

\textbf{Reflection Questions:}
\begin{enumerate}
\item When might you prefer mean vs. median as a measure of central tendency?
\item How can you verify that your statistical calculations are correct?
\item What's one new R function or concept you learned today?
\end{enumerate}

\textbf{Next Class Preview:}
We'll explore data visualization with histograms and boxplots, plus dive deeper into R programming basics.

\end{document}
